\begingroup
\newcommand{\IndexSurd}[1]{{}^{#1}\!\surd}
\clearpage\SourceLeaf{315}\NoSourceHead
\SourceGroup{II}
\HistoricTitle{[LÖSUNG DES UMKEHRPROBLEMS FÜR DAS ELLIPTISCHE\\ INTEGRAL ERSTER GATTUNG.]}
\HistoricCitation{[Aus Handbuch 16, Bb, Den astronomischen Wissenschaften gewidmet,\\ November 1801, S. 111—112]}
\par\indent Setzt man
\[\begin{gathered}
1+2x\cos\varphi+2x^4\cos2\varphi+2x^9\cos3\varphi+2x^{16}\cos4\varphi+\text{etc.}=P,\\[5pt]
\frac{dP}{d\varphi}=P',\qquad\frac{dP'}{d\varphi}=P'',
\end{gathered}\]
\noindent so wird
\[\begin{aligned}
&P'P'-PP''\\
&=(4xx+16x^8+36x^{18}+64x^{32}+\text{etc.})\\
&\qquad(1+2xx\cos2\varphi+2x^4\cos4\varphi+\text{etc.})\\
&\quad+(2x^{\frac12}+18x^{\frac92}+50x^{\frac{25}2}+98x^{\frac{49}2}+\quad{.}{.})\\
&\qquad(x^{\frac12}\cos\varphi+x^{\frac92}\cos3\varphi+x^{\frac{25}2}\cos5\varphi+\text{etc.}).
\end{aligned}\]
\noindent Ferner
\[\begin{aligned}
PP&=(1+2xx+2x^8+2x^{18}+\text{etc.})\\
&\qquad(1+2xx\cos2\varphi+2x^8\cos4\varphi+\text{etc.})\\
&\quad+(4x^{\frac12}+4x^{\frac92}+4x^{\frac{25}2}+4x^{\frac{49}2}+\text{etc.})\\
&\qquad(x^{\frac12}\cos\varphi+x^{\frac92}\cos3\varphi+x^{\frac{25}2}\cos5\varphi+\text{etc.}).
\end{aligned}\]
\noindent Noch
\[
2x^{\frac12}\cos\varphi+2x^{\frac92}\cos3\varphi+2x^{\frac{25}2}\cos5\varphi+2x^{\frac{49}2}\cos7\varphi+\text{etc.}=B
\]
\noindent gesetzt,
\[\begin{aligned}
BB={}&(2x+2x^9+2x^{25}+2x^{49}+\text{etc.})\\
&\qquad(1+2x^4\cos4\varphi+2x^{16}\cos8\varphi+2x^{36}\cos12\varphi+\text{etc.})\\
&+(1+2x^4+2x^{16}+2x^{36}+\text{etc.})\\
&\qquad(2x\cos2\varphi+2x^9\cos6\varphi+2x^{25}\cos10\varphi+\text{etc.})
\end{aligned}\]
\clearpage\SourceLeaf{316}\ArithmeticHead{309}{R}{LÖSUNG DES UMKEHRPROBLEMS.}
\noindent oder besser
\[\begin{aligned}
R={}&2x^{\frac14}\sin\tfrac12\varphi-2x^{\frac92}\sin\tfrac32\varphi\\
&+2x^{\frac{25}4}\sin\tfrac52\varphi-2x^{\frac{49}4}\sin\tfrac72\varphi+\text{etc.}
\end{aligned}\]
\noindent gesetzt, wird
\[\begin{aligned}
RR={}&(2x^{\frac12}+2x^{\frac92}+2x^{\frac{25}2}+2x^{\frac{49}2}+\text{etc.})\\
&\qquad(1+2xx\cos2\varphi+2x^8\cos4\varphi+\text{etc.})\\
&-(1+2xx+2x^8+2x^{18}+\text{etc.})\\
&\qquad(2x^{\frac12}\cos\varphi+2x^{\frac92}\cos3\varphi\\
&\qquad\quad+2x^{\frac{25}2}\cos5\varphi+2x^{\frac{49}2}\cos7\varphi+\text{etc.})
\end{aligned}\]
\noindent und
\[\begin{aligned}
&R'R'-RR''\\
&=\tfrac12(2x^{\frac12}+18x^{\frac92}+50x^{\frac{25}2}+98x^{\frac{49}2}+\text{etc.})\\
&\qquad(1+2xx\cos2\varphi+2x^8\cos4\varphi+2x^{18}\cos6\varphi+\text{etc.})\\
&\quad-(4xx+16x^8+36x^{18}+64x^{32}+\text{etc.})\\
&\qquad(2x^{\frac12}\cos\varphi+2x^{\frac92}\cos3\varphi+2x^{\frac{25}2}\cos5\varphi+\text{etc.}).
\end{aligned}\]
\par\indent Setzt man also
\[\begin{gathered}
1+2xx\cos2\varphi+2x^8\cos4\varphi+\text{etc.}=p,\qquad
1+2xx+2x^8+\text{etc.}=a,\\[5pt]
\frac{x\partial a}{\partial x}=a',\\[5pt]
2x^{\frac12}\cos\varphi+2x^{\frac92}\cos3\varphi+\text{etc.}=r,\qquad
2x^{\frac12}+2x^{\frac92}+\text{etc.}=b,\\[5pt]
\frac{x\partial b}{\partial x}=b',
\end{gathered}\]
\noindent so wird
\[\begin{aligned}
P'P'-PP''&=a'p+b'r,\qquad &PP&=ap+br,\\
R'R'-RR''&=b'p-a'r, &RR&=bp-ar,
\end{aligned}\]
\[\begin{aligned}
P^4+R^4&=(aa+bb)(pp+rr)\\
&=(aa+bb)^{\frac32}(1+2x\cos2\varphi+2x^4\cos4\varphi+2x^9\cos6\varphi+\text{etc.}),
\end{aligned}\]
\[
P'P'-PP''=\frac{aa'+bb'}{aa+bb}PP-\frac{ab'+ba'}{aa+bb}RR,
\]
\[
R'R'-RR''=\frac{ab'-ba'}{aa+bb}PP+\frac{aa'+bb'}{aa+bb}RR,
\]
\[\begin{aligned}
aa+bb&=(1+2x+2x^4+\text{etc.})^2,\\
aa'+bb'&=(1+2x+2x^4+\text{etc.})(2x+8x^4+\text{etc.}).
\end{aligned}\]
\clearpage\SourceLeaf{317}\ArithmeticHead{310}{L}{ANALYSIS.\quad NACHLASS.}
\par\indent Hier ist noch zu bemerken (wovon jedoch der Beweis tiefer liegt)
\[ab'-ba'=\tfrac12ab(a^4-b^4),\]
\noindent also
\[
\frac{P'P'-PP''}{PP}-\frac{R'R'-RR''}{RR}=-\frac12\left(\frac{PP}{RR}+\frac{RR}{PP}\right)ab(aa-bb).
\]
\noindent Noch findet man
\[\begin{aligned}
PR'-RP'={}&(x^{\frac18}+3x^{\frac98}-5x^{\frac{25}8}-7x^{\frac{49}8}+9x^{\frac{81}8}+{.}{.}{.})\\
&.\,(x^{\frac18}\cos\tfrac12\varphi-x^{\frac98}\cos\tfrac32\varphi\\
&\qquad-x^{\frac{25}8}\cos\tfrac52\varphi+x^{\frac{49}8}\cos\tfrac72\varphi+\text{etc.}).
\end{aligned}\]
\par\indent Das Quadrat des zweiten Factors im andern Theile der vorstehenden Gleichung wird
\[\begin{aligned}
&=\tfrac12(x^{\frac14}+x^{\frac94}+x^{\frac{25}4}+x^{\frac{49}4}+\text{etc.})\\
&\qquad(1-2x\cos2\varphi+2x^4\cos4\varphi-2x^9\cos6\varphi+\text{etc.})\\
&\quad+\tfrac12(1-2x+2x^4-2x^9+\text{etc.})\\
&\qquad(x^{\frac14}\cos\varphi+x^{\frac94}\cos3\varphi\\
&\qquad\quad+x^{\frac{25}4}\cos5\varphi+x^{\frac{49}4}\cos7\varphi[+\text{etc.}]).
\end{aligned}\]
\noindent Der erste Factor wird
\[=(aa+bb)\IndexSurd4\frac{(aa-bb)ab}{2}\quad?\]
\noindent Zusammen wird reductis reducendis
\[\begin{aligned}
(PR'-RP')^2&=\tfrac14(aa+bb)(ap-br)(bp+ar)\\
&=\tfrac14((aa-bb)PP+2abRR)(2abPP-(aa-bb)RR).
\end{aligned}\]
\par\indent Setzt man also
\[\frac RP\surd\frac{aa-bb}{2ab}=\sin\theta,\]
\noindent so wird
\[
\partial\varphi=\frac{2\partial\theta}{\surd\{(aa-bb)^2\cos\theta^2+(aa+bb)^2\sin\theta^2\}},
\]
\[
\cos\theta=\frac{2x^{\frac14}\cos\tfrac12\varphi+2x^{\frac94}\cos\tfrac32\varphi+2x^{\frac{25}4}\cos\tfrac52\varphi+\text{etc.}}{P}\surd\frac{aa+bb}{2ab}\mathpunct{\raisebox{.5ex}{.}}
\]
\ClosingRule
\clearpage
\endgroup
